\documentclass[a4paper,10pt]{article}
\usepackage[utf8]{inputenc}
\usepackage[margin=1.5cm, bottom=2cm]{geometry}
\usepackage{fancyhdr}
\usepackage[dvipsnames]{xcolor}
\usepackage{enumitem}
\usepackage{xcolor}
\usepackage[table]{xcolor}
\setlength{\parindent}{0pt}
\usepackage{makeidx}
\usepackage[most]{tcolorbox}
\newtcolorbox{osservazioni}{
  colback=gray!8,    % sfondo grigio chiaro
  colframe=white,      % nessun bordo
  arc=2mm,            % angoli smussati
  boxrule=0pt,        % spessore bordo 0
  fontupper=\footnotesize,
  left=2mm, right=2mm, top=2mm, bottom=2mm,
  before skip=7pt,   % niente spazio prima
  after skip=15pt     % niente spazio dopo
}
\newtcolorbox{osservazioni2}{
  colback=ForestGreen!4,    % sfondo grigio chiaro
  colframe=white,      % nessun bordo
  arc=2mm,            % angoli smussati
  boxrule=0pt,        % spessore bordo 0
  fontupper=\footnotesize,
  left=2mm, right=2mm, top=2mm, bottom=2mm,
  before skip=7pt,   % niente spazio prima
  after skip=15pt     % niente spazio dopo
}
\pagestyle{fancy}
\usepackage{tikz}
\usetikzlibrary{automata, positioning}
\pagestyle{fancy}

\usepackage{makeidx}

%%%%%%%%%%%%%%%%%%%
\newcommand{\EsercitazioneNum}{14}
\newcommand{\Data}{15/12/2025}
\newcommand{\DocType}{Esercitazione} % o "Eserciziario"
\newcommand{\FullTitle}{\DocType\ \EsercitazioneNum}

\newcommand{\Header}{
\noindent
  \small \textbf{Computabilità, Complessità e Logica - a.a. 2025/26} \hfill \textbf{\FullTitle} \\[0.2cm]
  \scriptsize Matteo Liotta, \texttt{matteo.liotta@studenti.units.it} \hfill \Data \\[0.2cm]
  \noindent\makebox[\linewidth]{\rule{\linewidth}{0.4pt}} \\[0.3cm]}

\newcommand{\NewPageHeader}{
  \vfill
  \noindent\makebox[\linewidth]{\rule{\linewidth}{0.4pt}}
  \newpage
  {\Header}
}
%%%%%%%%%%%%%%%%%%%

% Numero di pagina in basso al centro
\fancyhf{}
\fancyfoot[C] \thepage
\renewcommand{\footrulewidth}{0pt} % niente linea in basso
\renewcommand{\headrulewidth}{0pt} % niente linea in alto

\begin{document}

\footnotesize  % font generale più piccolo
\Header\\[0.7cm]


%%%%%%%%%%%%%%%% FOGLIO 2

    
    %% INDEX
    \addcontentsline{toc}{section}{Esercitazione 14: \textit{Universo di Herbrand e verifica della soddisfacibilità per LDP}}
    
    \noindent\Large \textbf{Esercitazione 14}: \textit{Simulazione d'Esame (appello 13/02/2024)} 
    
    {\normalsize
    \begin{itemize}[label=, leftmargin=0pt]
    
    %\item \textbf{\underline{Formule booleane cont'd.}}

    %% INDEX
    %\addcontentsline{toc}{subsection}{Esercizio 4.}
    %%
    %\item \textbf{\textcolor{black}{Esercizio 4.} (**)}\\
    %Si considerino le seguenti formule booleane:    
    %    $$\Big(((A \rightarrow (B \land \lnot C)) \lor (\overline A \lor B)) \rightarrow (\lnot A \lor C)\Big) \lor (B \land \lnot A) $$
    %    \noindent e
    %    $$\phi_2 = \big((\lnot A \lor B) \land (\lnot C \lor B)\big) \rightarrow ((C \lor B) \land A)$$
    
    
    %Si indichi, giustificando il motivo, se la prima è equivalente alla seconda.\\[0.2cm]

    \item \textbf{\underline{Totale: 33 punti}}\\[0.1cm]
    
    %% INDEX
    \addcontentsline{toc}{subsection}{Esercizio A1.} 
    %%
    \item \textbf{\textcolor{black}{Esercizio A1.} (6 punti)} \\ 	
    Nelle due coppie di seguito, si dica se la seconda formula è conseguenza logica della prima
    $$A\lor (B\rightarrow C) \quad \quad A\land B$$
    $$ (\lnot A\lor B\lor C) \land (A\lor B)\quad \quad B\lor C$$
    
    %% INDEX
    \addcontentsline{toc}{subsection}{Esercizio A2.} 
    %%
    \item \textbf{\textcolor{black}{Esercizio A2.} (7 punti)} \\ 	
    Si consideri il seguente insieme di clausole e si usi l'algoritmo di Davis-Puttnam descrivendo i passi in dettaglio per mostrare l'eventuale soddisfacibilità. Se così, si indichi una valutazione. Dire se l'insieme di clausole soddisfa un formato speciale.
    $$ S = \{\{\lnot P,Q\},\ \{\lnot P,R\}\},\ \{\lnot P,\lnot S\}\},\ \{R,S\},\ \{\lnot Q,\lnot R\},\ \{\lnot Q,S\}\}$$

    
    %% INDEX
    \addcontentsline{toc}{subsection}{Esercizio A3.} 
    %%
    \item \textbf{\textcolor{black}{Esercizio A3.} (8 punti)} \\ 	
    Si considerino gli alberi binari infiniti (ogni nodo ha un figlio destro e uno sinistro) etichettati su simboli proposizionali in $\Sigma$. Si consideri un simboli di predicato binario $A(x,y)$ che viene valutato a \textit{Vero} se $x$ e $y$ se sono valutati su due noti $n$ e $n'$ rispettivamente e $n'$ è un discendente (a distanza arbitraria positiva) dell'albero di $n$.\\
    Si introducano se necessario altri simbli di \textit{costante}, \textit{funzione} e \textit{predicato} -indicando chiaramente quale sia sia la loro interpretazione sul dominio-. 
    Si scrivano delle formule in logica dei predicati che esprimano le seguenti proprietà:
    \begin{enumerate}
        \item La radice e i suoi figli sono etichettati $P$ (in quei nodi è verificato il predicato $P$).
        \item C'è un sottoalbero i cui nodi sono tutti etichettati da $P$.
        \item C'è un cammino dalla radice con un numero infinito di nodi etichettati da $P$.
    \end{enumerate}
    
    %% INDEX
    \addcontentsline{toc}{subsection}{Esercizio A4.} 
    %%
    \item \textbf{\textcolor{black}{Esercizio A4.} (7 punti)} \\ 	
    Si consideri la seguente formula
    $$\lnot \bigg[\ \forall x.\bigg(P(x)\land (\ \forall y.Q(y)\rightarrow R(x,y)\ ) \land (\ \exists y. S(x,y)\rightarrow P(x)\ )\bigg)\ \bigg]$$
    \begin{enumerate}
        \item Si scriva la formula equivalente in forma prenessa.
        \item Si trasformi la formula prenessa in forma di clausole.
        \item Si scriva l'universo di Herbrand per la formula in forma di clausole
    \end{enumerate}

    %% INDEX
    \addcontentsline{toc}{subsection}{Esercizio A5.} 
    %%
    \item \textbf{\textcolor{black}{Esercizio A5.} (5 punti, \textit{facoltativo})} \\ 	
    Si consideri il seguente insieme di clausole:
    $$\{P(x),\ D(x)\},\quad \{\lnot P(x),\lnot D(x)\},\ \{\lnot P(x),\ D(succ(x))\},\quad \{\lnot D(x),\ P(succ(x))\}$$
    $$\{\lnot P(x),\ \lnot P(y),\ P(+(x,y))\}\quad \quad \{\lnot D(x),\ \lnot D(y),\ P(+(x,y))\}$$
    $$\{\lnot P(x),\ \lnot D(y),\ D(+(x,y))\}\quad\quad\{\lnot D(x),\ \lnot P(y),\ D(+(x,y))\}$$
    $$\{P(0)\}\quad \quad \{D(+(0,succ^2(0)\}$$
    Dove $x$ ed $y$ sono variabili, $0$ una costante (ad esempio lo zero dei naturali), $P$ e $D$ dei predicati, $succ$ e $+$ dei simboli di funzione. Si trovi una refutazione per l'insieme di clausole.

    
\end{itemize}
\vfill
\noindent\makebox[\linewidth]{\rule{\linewidth}{0.4pt}}
\end{document}